\section{Discussion}

\subsection{Main findings}

In this national registry study, we emulated 34 target trials of reperfusion treatment for ischaemic stroke, defined by stroke severity, type and timing of treatment, and use of perfusion imaging. Across most groups, treatment with IVT, MT or both was associated with more patients having a good outcome at discharge, lower discharge disability, fewer in-hospital deaths, shorter hospital stays and more discharges home.

We asked how these associations vary with stroke severity, with the type and timing of treatment, and across outcomes. First, they varied strongly with severity. They were small in mild stroke (NIHSS 1--5) and much larger in moderate to severe stroke. Second, they varied with treatment strategy. They were largest when IVT and MT were both given within 4.5 hours of onset, and were generally larger for early than for late treatment. Third, the pattern was broadly consistent across the five outcomes, with two exceptions: in-hospital death was slightly higher with IVT in mild stroke, and hospital stay was longer after late MT in the most severe strokes.

\subsection{Comparison with clinical trials}

The direction of our findings is consistent with trial evidence. Trials show that alteplase given within 4.5 hours improves outcome across a wide range of stroke severity\cite{emberson_effect_2014}, that thrombectomy substantially reduces disability in large-vessel occlusion\cite{goyal_endovascular_2016}, and that thrombectomy remains effective in later time windows in selected patients\cite{nogueira_thrombectomy_2018, albers_thrombectomy_2018}. We found associations in the same direction for each of these.

The size of our estimates is harder to compare with trials, and some estimates are larger than trial evidence would predict. In the pooled alteplase trials, a good outcome (mRS 0--1 at 3 to 6 months) was achieved by 32.9\% of treated patients and 23.1\% of controls when treatment was given within 3 hours, and by 35.3\% and 30.1\% when given between 3 and 4.5 hours\cite{emberson_effect_2014}. In our study, early IVT was associated with 16 to 20 percentage points more good outcomes in NIHSS 6--20. In the HERMES meta-analysis, 46.0\% of patients given thrombectomy and 26.5\% of controls were functionally independent at 90 days\cite{goyal_endovascular_2016}. Most patients in both arms also received IVT. We found differences of 23 to 34 percentage points for early IVT plus early MT, but compared with no reperfusion treatment at all, so a larger difference is expected. Our outcome definition and timing also differ from those used in trials: we measured outcome at discharge, and counted a return to pre-stroke mRS as a good outcome.

The clearest difference from trials is in mortality. Alteplase did not reduce 90-day mortality in the pooled trials (17.9\% with alteplase and 16.5\% with control)\cite{emberson_effect_2014}, and mortality did not differ significantly between arms in HERMES\cite{goyal_endovascular_2016}. In contrast, we found that in-hospital death was up to 12 percentage points lower with early IVT alone, and up to 21 percentage points lower with early IVT plus early MT. In-hospital death is not the same as 90-day mortality, but differences of this size are unlikely to be explained by treatment alone. They probably reflect, at least in part, the biases described below: patients who are dying or who are not for active treatment are less likely to receive reperfusion treatment. The mortality results in particular should therefore be interpreted with caution.

\subsection{Variation with stroke severity}

In mild stroke, the association with early IVT was small (3.0 percentage points more patients with a good outcome), and late IVT was not associated with benefit. About three quarters of untreated patients in this group had a good outcome, which leaves little room for improvement. In-hospital death was slightly higher with IVT, although deaths were rare in both arms. This is in keeping with the PRISMS trial, which found no benefit of alteplase in minor non-disabling stroke\cite{khatri_effect_2018}, and with the known risk of intracranial haemorrhage after alteplase\cite{emberson_effect_2014}. Two features of our analysis may also work against treatment in this group. Patients were matched on NIHSS band and not on exact score, and within the NIHSS 1--5 band treated patients had higher scores than their matched controls (mean 3.7 \textit{vs} 2.6 for early IVT; Supplementary Table~\ref{tab:supp_baseline_unmatched}). They are therefore likely to have had more disabling deficits. We also could not separate disabling from non-disabling deficits.

In moderate to severe stroke, treatment was associated with large differences in every outcome. The form of the association changed with severity. The difference in good outcome was generally greatest in NIHSS 11--20 and smaller in NIHSS 21--42, where a good outcome was uncommon in both arms. The differences in death and in discharge home or with community support were generally greatest in the most severe strokes.

\subsection{Variation with type and timing of treatment}

The combination of early IVT and early MT was associated with some of the largest improvements across outcomes. Within each severity band, associations were generally larger when treatment was given within 4.5 hours of onset than when it was given later or at an unknown time. This is consistent with trial evidence that the benefit of reperfusion treatment falls with time from onset\cite{emberson_effect_2014}.

These comparisons between trial groups are indirect and should be made with care. Each group was compared with its own untreated controls, and the patients who receive each treatment strategy differ. For example, between a quarter and a half of patients given MT without IVT were anticoagulated (Table~\ref{tab:baseline_matched}), which commonly rules out IVT, and patients given late MT may have been transferred from another hospital. Differences between groups may therefore reflect differences in patients as well as differences in treatment.

Associations were weaker, and often not statistically significant, in groups treated late with perfusion imaging. This should not be read as evidence that perfusion-selected treatment is ineffective. These groups were small, and patients who have perfusion imaging are likely to differ from other late-treated patients in ways we could not measure, such as unknown onset time, wake-up stroke or imaging findings. In addition, perfusion imaging is recorded in the thrombectomy section of SSNAP, so we could not identify patients selected for late IVT by perfusion imaging, as in EXTEND\cite{ma_thrombolysis_2019, campbell_extending_2019}.

\subsection{Outcomes beyond functional status}

This study reports outcomes that are rarely available from trials but matter to patients and health services. Treatment was associated with large shifts in discharge destination, with more patients going home or with community team support, and fewer going to a care home or dying in hospital. Hospital stays were generally shorter, with reductions of up to about 18 days with early IVT plus early MT. The exception was late MT in the most severe strokes, which was associated with longer stays. This group also had fewer in-hospital deaths, so the longer stay may reflect more severely affected patients surviving and needing rehabilitation. Length of stay therefore needs to be interpreted alongside mortality.

\subsection{Strengths and limitations}

The main strengths of this study are its use of national data covering most stroke admissions in England and Wales, the explicit specification of each target trial, and the reporting of several outcomes for many treatment and severity groups. Almost all treated patients were matched, so the results apply to nearly the whole treated population and not to a selected subset.

The main limitation is residual confounding. Patients were matched on a small number of variables: pre-stroke disability, NIHSS band, age band, AF diagnosis and anticoagulant use for AF. We did not adjust for other comorbidities or for process measures. Patients who were not treated may differ from treated patients in ways not captured, such as arrival time, contraindications, large-vessel occlusion, imaging findings or the reason treatment was not given. In particular, untreated controls for MT groups were not required to have a large-vessel occlusion, so their prognosis may differ from that of treated patients with the same NIHSS. Some untreated patients will not have been treated because they were judged too unwell to benefit, which would bias results in favour of treatment.

Matching on NIHSS band and not on exact score left imbalance in severity within some bands (Supplementary Table~\ref{tab:supp_baseline_unmatched}). In NIHSS 1--5, treated patients had higher scores than controls, which works against treatment. In NIHSS 21--42, treated patients had lower scores than controls in every group (mean about 24 \textit{vs} 25), which favours treatment and may explain part of the large differences seen in the most severe strokes, including in mortality. Diabetes was also less common in treated patients in most groups. Age is recorded only in 5-year bands.

Time zero was stroke onset, but treatment status was defined from treatment received at any point during the admission, and not from a decision made at onset. Patients who died or deteriorated very early, before treatment could be given, are more likely to be in the untreated group. This can make treatment look more beneficial than it is, a form of immortal time bias. Patients treated when the time of stroke onset was not known were grouped with late treatment. The late groups therefore mix treatment known to be late with treatment of unknown timing, some of which may have been given early.

Outcomes were measured at discharge and not at 90 days, as in most trials. Discharge outcomes depend on local discharge practice and the availability of rehabilitation and community services, and they may not reflect longer-term recovery. Finally, the same untreated patient could serve as a control in more than one trial group, so the groups are not independent, and many comparisons were made. We used Holm adjustment within each outcome, but small groups had limited power and wide confidence intervals.

\subsection{Implications and further work}

For clinicians and patients, these results describe what happened in routine care to patients like those who were treated, by stroke severity and treatment strategy, across outcomes that trials seldom report. They are consistent with current practice of giving reperfusion treatment as early as possible, and suggest that the expected gain from IVT in mild stroke is small. For service planners, the associations with length of stay and discharge destination are relevant to planning hospital beds, rehabilitation and community services as access to thrombectomy grows.

Further work should test how robust these estimates are. Useful analyses would include matching on exact NIHSS score, adjustment for more patient characteristics, restriction of controls for MT groups to patients with a recorded large-vessel occlusion, restriction to patients with a known onset time, and methods that handle the timing of treatment explicitly to reduce immortal time bias. Longer-term outcomes would show whether the differences seen at discharge persist.

\subsection{Conclusions}

In routine care in England and Wales, reperfusion treatment for ischaemic stroke was associated with better outcomes at discharge across most combinations of stroke severity, treatment type and timing. Associations were small in mild stroke and greatest when IVT and MT were combined early in moderate to severe stroke. By reporting death, disability, length of stay and discharge destination for each group, these results add detail that trials do not provide and that may help clinicians, patients and service planners. Because these are observational estimates that may be affected by residual confounding, they should be read as associations that complement, and do not replace, randomised evidence.
